% ==============================================================
\section{The Finite-Horizon Problem}
\label{sec:ac-T-period}
% ==============================================================

We now extend the two-period analysis to a finite horizon of $T+1$ periods, $t \in \{0, 1, \ldots, T\}$. This extension reveals the full dynamic structure of adjustment costs: how investment at any date responds to productivity changes at any other date, and how Tobin's $Q$ evolves along the transition path.

% --------------------------------------------------------------
\subsection{Setup}
\label{subsec:ac-T-setup}
% --------------------------------------------------------------

\paragraph{Technology.}
Production is linear in capital with productivity $A_t$:
\begin{equation}
Y_t = A_t k_t.
\label{eq:ac-T-production}
\end{equation}
Capital fully depreciates ($\delta = 1$), so next-period capital equals current investment:
\begin{equation}
k_{t+1} = I_t, \qquad t = 0, 1, \ldots, T-1.
\label{eq:ac-T-capital}
\end{equation}
The CES aggregator governs the transformation of output in each period:
\begin{equation}
\Hfun(C_t, I_t) = Y_t, \qquad t = 0, 1, \ldots, T-1.
\label{eq:ac-T-resource}
\end{equation}
In the terminal period $T$, there is no investment: $C_T = Y_T$.

\paragraph{Preferences.}
The household has time-separable preferences with discount factor $\beta \in (0,1)$ and period utility $u(C) = \frac{C^{1-1/\theta} - 1}{1 - 1/\theta}$:
\begin{equation}
U = \sum_{t=0}^{T} \beta^t \, u(C_t).
\label{eq:ac-T-utility}
\end{equation}

\paragraph{State and Control Variables.}
The state variable is capital $k_t$. The control variables are consumption $C_t$ and investment $I_t$, linked by the CES constraint. Given the constraint, we can take $I_t$ as the single control: $C_t$ is then determined by $\Hfun(C_t, I_t) = A_t k_t$.

% --------------------------------------------------------------
\subsection{The Recursive Formulation}
\label{subsec:ac-T-recursive}
% --------------------------------------------------------------

\paragraph{The Bellman Equation.}
Define $V_t(k)$ as the value of entering period $t$ with capital $k$. The Bellman equation is:
\begin{equation}
V_t(k) = \max_{I} \left\{ u(C) + \beta V_{t+1}(I) \right\} \quad \text{subject to} \quad \Hfun(C, I) = A_t k,
\label{eq:ac-T-bellman}
\end{equation}
with terminal condition $V_T(k) = u(A_T k)$.

\paragraph{Eliminating $C$.}
From the CES constraint, consumption is implicitly defined by $\Hfun(C, I) = Y$. For the canonical CES:
\begin{equation}
C = \left( \frac{Y^\psi - \omega_I^{1-\psi} I^\psi}{\omega_C^{1-\psi}} \right)^{1/\psi} \equiv C(I; Y).
\label{eq:ac-T-C-implicit}
\end{equation}

The Bellman equation becomes:
\begin{equation}
V_t(k) = \max_{I} \left\{ u\bigl(C(I; A_t k)\bigr) + \beta V_{t+1}(I) \right\}.
\label{eq:ac-T-bellman-reduced}
\end{equation}

\paragraph{First-Order Condition.}
Differentiating with respect to $I$:
\begin{equation}
u'(C) \frac{\partial C}{\partial I} + \beta V_{t+1}'(I) = 0.
\label{eq:ac-T-foc}
\end{equation}

From the implicit function theorem applied to $\Hfun(C, I) = Y$:
\begin{equation}
\frac{\partial C}{\partial I} = -\frac{\partial \Hfun / \partial I}{\partial \Hfun / \partial C} = -\frac{1}{Q},
\label{eq:ac-T-dCdI}
\end{equation}
where $Q$ is Tobin's $Q$ as defined in \cref{eq:ac-tobin-q}. The first-order condition becomes:
\begin{equation}
\frac{u'(C_t)}{Q_t} = \beta V_{t+1}'(k_{t+1}).
\label{eq:ac-T-foc-Q}
\end{equation}

\paragraph{Envelope Condition.}
Differentiating the value function with respect to $k$:
\begin{equation}
V_t'(k) = u'(C) \frac{\partial C}{\partial Y} \cdot A_t = u'(C_t) \cdot \frac{A_t}{Q_t} \cdot \frac{1}{1 - s_I},
\label{eq:ac-T-envelope-raw}
\end{equation}
where $s_I = \omega_I^{1-\psi} I^\psi / Y^\psi$ is the investment share in the CES aggregator.

For the CES function, $\partial C / \partial Y = (C/Y) / (1 - s_I)$ by homogeneity arguments. At the optimum, the envelope condition simplifies. Using the homogeneity of $\Hfun$:
\begin{equation}
V_t'(k) = \frac{u'(C_t)}{Q_t} \cdot A_t.
\label{eq:ac-T-envelope}
\end{equation}

% --------------------------------------------------------------
\subsection{The Euler Equation and Tobin's $Q$}
\label{subsec:ac-T-euler}
% --------------------------------------------------------------

Combining the first-order condition \cref{eq:ac-T-foc-Q} with the envelope condition \cref{eq:ac-T-envelope} (shifted forward one period):
\begin{equation}
\frac{u'(C_t)}{Q_t} = \beta \cdot \frac{u'(C_{t+1})}{Q_{t+1}} \cdot A_{t+1}.
\label{eq:ac-T-euler-combined}
\end{equation}

Rearranging:
\begin{equation}
\boxed{u'(C_t) = \beta \frac{A_{t+1}}{Q_{t+1}} Q_t \, u'(C_{t+1}).}
\label{eq:ac-T-euler}
\end{equation}

\paragraph{Interpretation.}
The Euler equation has a natural interpretation in terms of returns:
\begin{itemize}[nosep]
\item $Q_t$ is the cost of acquiring one unit of capital at date $t$ (in consumption units).
\item $A_{t+1}$ is the marginal product of that capital at date $t+1$.
\item $Q_{t+1}$ is the cost of capital at date $t+1$---relevant if the household could instead acquire capital next period.
\end{itemize}

Define the \emph{return on capital}:
\begin{equation}
R_{t+1} \equiv \frac{A_{t+1}}{Q_{t+1}} \cdot Q_t.
\label{eq:ac-T-return}
\end{equation}

The Euler equation becomes the standard form:
\begin{equation}
u'(C_t) = \beta R_{t+1} \, u'(C_{t+1}).
\label{eq:ac-T-euler-standard}
\end{equation}

\begin{calloutimportant}[The Return with Adjustment Costs]
The return $R_{t+1} = \frac{A_{t+1}}{Q_{t+1}} Q_t$ has two components:
\begin{enumerate}[nosep]
\item \textbf{Dividend yield}: $A_{t+1}/Q_{t+1}$ is the marginal product per unit of capital value.
\item \textbf{Capital gain}: $Q_t/Q_{t+1}$ captures changes in Tobin's $Q$. If $Q$ is falling ($Q_{t+1} < Q_t$), the return is higher---capital acquired cheaply today can be ``sold'' at a higher implicit price tomorrow.
\end{enumerate}

In the frictionless case ($\eta = \infty$), $Q_t = Q_{t+1} = 1$ and $R_{t+1} = A_{t+1}$. With adjustment costs, returns include capital gains/losses on $Q$.
\end{calloutimportant}

% --------------------------------------------------------------
\subsection{The Sequence Formulation}
\label{subsec:ac-T-sequence}
% --------------------------------------------------------------

An alternative to the recursive approach is to solve the problem in sequence space. This formulation connects directly to the Jacobian methods discussed in earlier chapters.

\paragraph{The Planner's Problem.}
The planner maximizes:
\begin{equation}
\max_{\{C_t, I_t, k_{t+1}\}} \sum_{t=0}^{T} \beta^t u(C_t)
\label{eq:ac-T-planner-obj}
\end{equation}
subject to:
\begin{align}
\Hfun(C_t, I_t) &= A_t k_t, \qquad t = 0, \ldots, T-1, \label{eq:ac-T-planner-ces} \\
k_{t+1} &= I_t, \qquad t = 0, \ldots, T-1, \label{eq:ac-T-planner-accum} \\
C_T &= A_T k_T, \label{eq:ac-T-planner-terminal}
\end{align}
with $k_0$ given.

\paragraph{Lagrangian.}
Form the Lagrangian with multipliers $\lambda_t$ on the CES constraints and $\mu_t$ on the capital accumulation equations:
\begin{equation}
\mathcal{L} = \sum_{t=0}^{T} \beta^t u(C_t) + \sum_{t=0}^{T-1} \lambda_t \left[ A_t k_t - \Hfun(C_t, I_t) \right] + \sum_{t=0}^{T-1} \mu_t \left[ I_t - k_{t+1} \right].
\label{eq:ac-T-lagrangian}
\end{equation}

\paragraph{First-Order Conditions.}
\begin{align}
C_t: & \quad \beta^t u'(C_t) = \lambda_t \frac{\partial \Hfun}{\partial C_t}, \qquad t = 0, \ldots, T-1, \label{eq:ac-T-foc-C} \\
I_t: & \quad \lambda_t \frac{\partial \Hfun}{\partial I_t} = \mu_t, \qquad t = 0, \ldots, T-1, \label{eq:ac-T-foc-I} \\
k_{t+1}: & \quad \mu_t = \lambda_{t+1} A_{t+1}, \qquad t = 0, \ldots, T-2, \label{eq:ac-T-foc-k} \\
C_T: & \quad \beta^T u'(C_T) = \mu_{T-1} A_T. \label{eq:ac-T-foc-CT}
\end{align}

\paragraph{Recovering Tobin's $Q$.}
From \cref{eq:ac-T-foc-C,eq:ac-T-foc-I}:
\begin{equation}
\frac{\mu_t}{\lambda_t} = \frac{\partial \Hfun / \partial I_t}{\partial \Hfun / \partial C_t} \cdot \frac{\beta^t u'(C_t)}{\beta^t u'(C_t)} = \frac{1}{Q_t}.
\label{eq:ac-T-mu-lambda}
\end{equation}

The ratio $\mu_t / \lambda_t = 1/Q_t$ is the shadow price of investment relative to the shadow price of output. Equivalently, $\mu_t$ is the shadow value of capital (in utils), and $\lambda_t Q_t$ is the shadow value of consumption.

\paragraph{The Euler Equation (Again).}
From \cref{eq:ac-T-foc-k}: $\mu_t = \lambda_{t+1} A_{t+1}$. Substituting $\mu_t = \lambda_t / Q_t$:
\begin{equation}
\frac{\lambda_t}{Q_t} = \lambda_{t+1} A_{t+1}.
\label{eq:ac-T-lambda-recursion}
\end{equation}

From \cref{eq:ac-T-foc-C}: $\lambda_t = \beta^t u'(C_t) / (\partial \Hfun / \partial C_t)$. Using $\partial \Hfun / \partial C = 1/Q$ (up to scaling):
\begin{equation}
\lambda_t \propto \beta^t u'(C_t) Q_t.
\label{eq:ac-T-lambda-prop}
\end{equation}

Substituting into \cref{eq:ac-T-lambda-recursion} and simplifying yields the Euler equation \cref{eq:ac-T-euler}.

% --------------------------------------------------------------
\subsection{Homothetic Structure and the Normalized Problem}
\label{subsec:ac-T-homothetic}
% --------------------------------------------------------------

With homothetic preferences (CRRA) and a homogeneous-of-degree-one production/transformation technology, the problem admits a normalized formulation.

\paragraph{Guess for the Value Function.}
Conjecture that the value function takes the form:
\begin{equation}
V_t(k) = v_t \cdot \frac{k^{1-1/\theta}}{1 - 1/\theta},
\label{eq:ac-T-value-guess}
\end{equation}
where $v_t$ is a time-dependent coefficient to be determined.

\paragraph{Normalized Variables.}
Define the consumption-output ratio $\bar{c}_t \equiv C_t / Y_t$ and the investment-output ratio $\bar{\imath}_t \equiv I_t / Y_t$. The CES constraint becomes:
\begin{equation}
\Hfun(\bar{c}_t, \bar{\imath}_t) = 1,
\label{eq:ac-T-normalized-ces}
\end{equation}
which implicitly defines $\bar{c}_t$ as a function of $\bar{\imath}_t$.

\paragraph{Tobin's $Q$ in Normalized Form.}
\begin{equation}
Q_t = \frac{\omega_I^{1-\psi}}{\omega_C^{1-\psi}} \left( \frac{\bar{c}_t}{\bar{\imath}_t} \right)^{1-\psi}.
\label{eq:ac-T-Q-normalized}
\end{equation}

\paragraph{The Normalized Bellman Equation.}
Substituting the guess into the Bellman equation and using $k_{t+1} = I_t = \bar{\imath}_t Y_t = \bar{\imath}_t A_t k_t$:
\begin{equation}
v_t \frac{k^{1-1/\theta}}{1-1/\theta} = \max_{\bar{\imath}} \left\{ \frac{(\bar{c} \cdot A_t k)^{1-1/\theta}}{1-1/\theta} + \beta v_{t+1} \frac{(\bar{\imath} \cdot A_t k)^{1-1/\theta}}{1-1/\theta} \right\}.
\label{eq:ac-T-bellman-sub}
\end{equation}

Factoring out $k^{1-1/\theta}$ and $A_t^{1-1/\theta}$:
\begin{equation}
v_t = A_t^{1-1/\theta} \max_{\bar{\imath}} \left\{ \frac{\bar{c}^{1-1/\theta} + \beta v_{t+1} \bar{\imath}^{1-1/\theta}}{1-1/\theta} \right\} \cdot (1-1/\theta).
\label{eq:ac-T-bellman-normalized}
\end{equation}

This confirms the guess: the value function is linear in $k^{1-1/\theta}$, and the coefficients $v_t$ satisfy a recursion that depends only on $A_t$ and $\beta$, not on $k$.

\paragraph{The Normalized Problem.}
The coefficients $\{v_t\}$ solve a finite-dimensional problem:
\begin{equation}
v_t = \max_{\bar{\imath}_t} \left\{ \bar{c}_t^{1-1/\theta} + \beta v_{t+1} (\bar{\imath}_t A_t)^{1-1/\theta} \right\} \quad \text{s.t.} \quad \Hfun(\bar{c}_t, \bar{\imath}_t) = 1,
\label{eq:ac-T-v-recursion}
\end{equation}
with terminal condition $v_T = A_T^{1-1/\theta}$.

\begin{calloutnote}[Dimensionality Reduction]
The homothetic structure reduces the state space from $(k_t, t)$ to just $t$. The normalized investment rate $\bar{\imath}_t$ depends only on time (through $A_t$, $\beta$, and the terminal condition), not on the level of capital. This is the dynamic analog of the ``linearity in wealth'' property in consumption-savings problems.

This reduction is valuable computationally: instead of solving a value function over a continuous state space, we solve a sequence of static optimization problems, one for each period.
\end{calloutnote}

% --------------------------------------------------------------
\subsection{Steady State and Dynamics}
\label{subsec:ac-T-dynamics}
% --------------------------------------------------------------

Consider the case where productivity is constant: $A_t = A$ for all $t$. In an infinite-horizon version (or for $T$ large), the economy converges to a steady state.

\paragraph{Steady-State Conditions.}
In steady state, $\bar{c}_t = \bar{c}$, $\bar{\imath}_t = \bar{\imath}$, $Q_t = Q$, and $v_t = v$ for all $t$. The Euler equation \cref{eq:ac-T-euler} becomes:
\begin{equation}
1 = \beta \frac{A}{Q}.
\label{eq:ac-T-ss-euler}
\end{equation}

Thus, in steady state:
\begin{equation}
Q^{ss} = \beta A.
\label{eq:ac-T-Q-ss}
\end{equation}

Tobin's $Q$ equals the discounted marginal product of capital. If $\beta A > 1$, steady-state $Q > 1$: investment goods are expensive relative to the frictionless benchmark.

\paragraph{Steady-State Investment Rate.}
From the definition of $Q$:
\begin{equation}
\beta A = \frac{\omega_I^{1-\psi}}{\omega_C^{1-\psi}} \left( \frac{\bar{c}^{ss}}{\bar{\imath}^{ss}} \right)^{1-\psi}.
\label{eq:ac-T-Q-ss-def}
\end{equation}

Combined with the CES constraint $\Hfun(\bar{c}^{ss}, \bar{\imath}^{ss}) = 1$, this pins down the steady-state allocation.

\paragraph{Transitional Dynamics.}
Away from steady state, the system exhibits transitional dynamics. Starting from an initial capital stock $k_0 \neq k^{ss}$:
\begin{itemize}[nosep]
\item If $k_0 < k^{ss}$: The economy invests heavily, driving up $Q$. High $Q$ means low returns $A/Q$, which moderates investment over time.
\item If $k_0 > k^{ss}$: The economy disinvests (if possible) or lets capital depreciate, with low investment rates and low $Q$.
\end{itemize}

The adjustment cost (finite $\eta$) slows the transition: the economy cannot jump immediately to steady state because high investment rates make investment expensive at the margin.

% --------------------------------------------------------------
\subsection{Summary}
\label{subsec:ac-T-summary}
% --------------------------------------------------------------

\begin{table}[H]
\centering
\caption{The Finite-Horizon Problem with CES Adjustment Costs}
\label{tab:ac-T-summary}
\renewcommand{\arraystretch}{1.6}
\begin{tabular}{@{}ll@{}}
\toprule
\textbf{Object} & \textbf{Formula} \\
\midrule
Bellman equation & $V_t(k) = \max_{I} \left\{ u(C) + \beta V_{t+1}(I) \right\}$ s.t. $\Hfun(C,I) = A_t k$ \\[6pt]
Euler equation & $u'(C_t) = \beta \dfrac{A_{t+1}}{Q_{t+1}} Q_t \, u'(C_{t+1})$ \\[10pt]
Return on capital & $R_{t+1} = \dfrac{A_{t+1}}{Q_{t+1}} Q_t$ \\[10pt]
Steady-state $Q$ & $Q^{ss} = \beta A$ \\[6pt]
Value function form & $V_t(k) = v_t \cdot \dfrac{k^{1-1/\theta}}{1-1/\theta}$ \\[8pt]
\bottomrule
\end{tabular}
\end{table}

\begin{callouttip}[Key Insights]
The finite-horizon extension reveals:
\begin{itemize}[nosep]
\item \textbf{Returns include capital gains}: The return $R_{t+1} = (A_{t+1}/Q_{t+1}) Q_t$ has both a dividend yield and a capital gain component from changes in $Q$.
\item \textbf{Homothetic reduction}: With CRRA utility and CES transformation, the value function is homogeneous in capital, reducing the problem to a sequence of normalized static problems.
\item \textbf{Steady-state $Q$}: In steady state, $Q = \beta A$---Tobin's $Q$ equals the discounted marginal product.
\item \textbf{Gradual adjustment}: Finite transformation elasticity $\eta$ implies gradual convergence to steady state; the economy cannot jump instantaneously.
\end{itemize}
\end{callouttip}
